\begin{abstract}
GPU workloads with large memory footprints frequently suffer from redundant L2 TLB misses
in which a recently evicted translation is immediately re-walked at full page-walk cost.
We characterize these \emph{dead-entry misses} across 24 GPU workloads, finding they
account for up to 99\% of L2 TLB misses in the most TLB-sensitive applications, yet their
performance impact varies widely depending on memory access structure.
Workloads where warps share the same virtual page suffer from burst amplification, where
a single eviction stalls many warps simultaneously waiting for one translation to return.
In contrast, workloads where each warp accesses a distinct set of pages face a
capacity-overflow problem that no replacement policy can resolve, a distinction validated
by huge page experiments.
Building on this two-class taxonomy, we design \emph{DEPOT} (Dead-Entry PrOTection),
a 1\,KB Bloom filter mechanism that prevents recently evicted translations from being
displaced immediately upon reinstallation, delivering up to 72\% IPC improvement on
interference-driven workloads with zero overhead on others, and composing with the state-of-the-art TLB prefetching and compaction mechanism,
for 2 to 7\% additional gain.
\end{abstract}
